% Ejemplo con FECHAS: parcelación de Filosofía 10.°, periodo I, con texto de relleno.
% \fechainicio pone el rango de cada semana; \receso ocupa una semana sin clase.
\documentclass[parcelacion]{planeacion}
\usepackage{lipsum}
\materia{filosofia}
\grado{10.°} \periodo{I} \anio{2027}
\newcommand{\lorem}[2][1]{\lipsum[#1][#2]}
\fechainicio{2027-01-25}
\nota{\lipsum[30][1-2]}
\begin{document}
\begin{parcelacion}
  \semana{1}{\textbf{Lorem ipsum.} \lorem[3]{1}}{\lorem[4]{1-2}}{\lorem[5]{1}}{\lorem[6]{1}}
  \semana{2}{\textbf{Dolor sit amet.} \lorem[7]{1}}{\lorem[8]{1-2}}{\lorem[9]{1}}{\lorem[10]{1}}
  \semana[q]{3}{\textbf{Consectetuer.} \lorem[11]{1}}{\lorem[12]{1-2}}{\lorem[13]{1}}{\lorem[14]{1}}
  \semana{4}{\textbf{Adipiscing elit.} \lorem[15]{1}}{\lorem[16]{1-2}}{\lorem[17]{1}}{\lorem[18]{1}}
  \semana[i]{5}{\textbf{Ut purus elit.} \lorem[19]{1}}{\lorem[20]{1-2}}{\lorem[21]{1}}{\lorem[22]{1}}
  \semana[t]{6}{\textbf{Vestibulum.} \lorem[23]{1}}{\lorem[24]{1-2}}{\lorem[25]{1}}{\lorem[26]{1}}
  \semana{7}{\textbf{Ut placerat.} \lorem[27]{1}}{\lorem[28]{1-2}}{\lorem[29]{1}}{\lorem[31]{1}}
  \semana[p]{8}{\textbf{Curabitur.} \lorem[32]{1}}{\lorem[33]{1-2}}{\lorem[34]{1}}{\lorem[35]{1}}
  \receso{RECESO ESCOLAR — SEMANA SANTA}
  \semana[e]{9}{\textbf{Nulla malesuada.} \lorem[36]{1}}{\lorem[37]{1-2}}{\lorem[38]{1}}{\lorem[39]{1}}
  \semana[r]{10}{\textbf{Evaluación.} \lorem[40]{1}}{\lorem[41]{1-2}}{\lorem[42]{1}}{\lorem[43]{1}}
\end{parcelacion}
\end{document}
